% =============================================================================
% Resumo
% =============================================================================

\section*{Resumo}
\addcontentsline{toc}{section}{Resumo}

Este relatório apresenta o andamento parcial do estudo experimental que compara soluções de bancos de dados vetoriais para operações de busca semântica, com foco em aplicações de Inteligência Artificial baseadas em embeddings e \textit{Large Language Models} (LLMs). São avaliados um banco relacional com extensão vetorial (PostgreSQL + pgvector) e bancos vetoriais especializados (Qdrant e Weaviate), em cenários típicos de sistemas de \textit{Retrieval-Augmented Generation} (RAG) e busca semântica em documentos. A comparação considera métricas de latência, \textit{throughput}, qualidade da recuperação (recall@K) e consumo de recursos. Esta entrega parcial cobre a revisão bibliográfica e o planejamento, a preparação de dados e do ambiente experimental reproduzível e a implementação dos \textit{scripts} de \textit{benchmark} dos cenários de busca pura e de busca com filtros. A execução dos experimentos em escala, os experimentos de carga mista e a análise comparativa final serão apresentados no relatório final, em dezembro de 2026.

\noindent\textbf{Palavras-chave:} bancos de dados vetoriais; busca semântica; pgvector; Qdrant; Weaviate; RAG; \textit{benchmarking}.
